% Replacement prose for EXP-2; source V4 remains unchanged.
Forward evaluation of the same 1,842 unique feasible designs shows that class H attains the maximum generalized-cost ratio in every OD--period cell and at every design's worst-case cell. All 341 lower-left designs reduce low-income trips. Holding the class participation weights fixed at their same-state P0 values prevents the cost ratio from falling in 17 of these designs (5.0%), demonstrating a composition effect in that subset. In the remaining 324, the cost ratio still falls at fixed weights, so low-income exit alone does not explain the decline. These are accounting comparisons within the evaluated library, rather than causal estimates.

Using a common reference in-vehicle VOT of 97.6 CNY/h for all classes, together with common population-weighted reference valuations for other time attributes, the L-specific generalized-cost ratio falls while L participation falls in 660 designs. The corresponding counts are 728 for L's mean monetary cost per realized trip and 327 for the fixed-weight version of the original ratio. The diagnostics therefore support both high-VOT-class dominance of the original index and a participation-composition effect in a smaller subset; monitoring L participation remains necessary alongside realized-traveler burden measures.
